% Working proof-authoring fragment, not an adopted handbook input.
% Owner: lambda_hol_metatheory_author; 2026-10-02.
% Exact source: Barendregt--Geuvers 2001, Figure 2 and §§3.2--3.4.
% No new calculus rules are introduced here. Paragraph names are working IDs.

\section*{Working construction: the exact three-sort lambdaHOL metatheory}

This fragment treats only the beta-only Figure~2 system. Its sorts are
\(\mathrm{Prop},\mathrm{Type},\mathrm{Type}'\), its axioms are
\(\mathrm{Prop}:\mathrm{Type}\), \(\mathrm{Type}:\mathrm{Type}'\),
and its product triples are exactly
\((\mathrm{Type},\mathrm{Type},\mathrm{Type})\),
\((\mathrm{Prop},\mathrm{Prop},\mathrm{Prop})\), and
\((\mathrm{Type},\mathrm{Prop},\mathrm{Prop})\).
There are no definitions, delta contractions, recursors, or equality
reflection. These constructions do not assert a result for an arbitrary
PTS or for the informal beta/iota/delta regime of source section~2.
The raw annotated pseudo-term confluence and entire-context substitution
proofs are separately supplied by the source author; they are used, not
reproved below. None of W1--W7 assumes normalization or checker completeness.

\paragraph{W1. Valid contexts, projection, and thinning.}
Define a valid context inductively: the empty context is valid; a fresh
extension \(\Gamma,x:A\) is valid when \(\Gamma\) is valid and
\(\Gamma\vdash A:s\) for a source sort \(s\). Every context in a
source derivation is valid. This follows by induction on that derivation:
the sort cases have empty context; variable and weakening exhibit the
annotation certificate; product and abstraction inherit their base
context from a premise; application and conversion inherit their common
context. Conversely every valid context admits
\(\Gamma\vdash\mathrm{Prop}:\mathrm{Type}\): start from the axiom and
weaken through its annotation certificates. Thus this explicit validity
is equivalent to the source's occurrence-in-a-judgment definition.

If \(x:A\) occurs in a valid context, the variable rule types \(x:A\)
at its own prefix and successive weakenings carry that judgment through
the suffix. This proves projection without a new primitive rule.
If \(\Gamma\) and \(\Delta\) are valid, every declaration of
\(\Gamma\) occurs literally in \(\Delta\), and
\(\Gamma\vdash M:A\), then \(\Delta\vdash M:A\).
Induct on the displayed derivation. For a sort, weaken its empty-context
axiom through \(\Delta\). For a variable, use the projection just proved.
A weakening case uses the induction hypothesis on its subject premise;
its formerly added declaration is already present in \(\Delta\).
Product and abstraction use the induction hypotheses on their premises,
renaming the bound variable fresh for \(\Delta\); the new annotation
judgment makes \(\Delta,x:B\) valid, so the body induction applies
there. Application combines the two induction judgments. Conversion
combines its subject and well-typed result-type induction judgments and
the unchanged raw beta equality. This covers all Figure~2 rules.
In particular this thinning proof does not infer strengthening by
substituting a possibly nonexistent inhabitant.

\paragraph{W2. Generation and product injectivity.}
For every source judgment, generation has the following exact content.
For a variable \(x:T\) in the context, any type of \(x\) is
beta-equal to \(T\). Any type of \(\mathrm{Prop}\) is beta-equal to
\(\mathrm{Type}\); any type of \(\mathrm{Type}\) is beta-equal to
\(\mathrm{Type}'\); no judgment has subject \(\mathrm{Type}'\).
For a product \(\Pi x:B.C\), there are premises
\(\Gamma\vdash B:s_1\),
\(\Gamma,x:B\vdash C:s_2\), an allowed pair, and its conclusion type
is beta-equal to \(s_2\). For an abstraction \(\lambda x:B.P\),
there is a \(C\) with \(\Gamma,x:B\vdash P:C\),
\(\Gamma\vdash\Pi x:B.C:s\), and conclusion type beta-equal to
\(\Pi x:B.C\). For an application \(PQ\), there are
\(\Gamma\vdash P:\Pi x:B.C\), \(\Gamma\vdash Q:B\), and its
conclusion type is beta-equal to \(C[Q/x]\).
The displayed constructor judgment at this generated type is itself
derivable in \(\Gamma\).

Prove these statements by induction on the typing derivation. Its
constructor rule supplies the witnesses. A final conversion composes
the generated type equality with its own equality. A final weakening
carries every generated premise into the enlarged valid context by W1,
including the fresh binder extension, and leaves the equalities intact.
No other rule has the selected subject constructor. These are all cases.

Raw beta confluence gives product injectivity without normalization.
If \(\Pi x:B.C=_\beta\Pi x:D.E\), choose a common reduct.
A product has no root beta contraction, and every reduct still has
the product constructor with a reduced domain and body. After choosing
the same fresh binder, the common reduct therefore identifies common
reducts of \(B,D\) and of \(C,E\). Hence
\(B=_\beta D\), \(C=_\beta E\). Distinct sorts are not beta-equal,
and a product is not beta-equal to a sort, by the same argument.

\paragraph{W3. Type correctness.}
If \(\Gamma\vdash M:A\), then either \(A\) is a sort or
\(\Gamma\vdash A:s\) for a sort \(s\).
Induct on the derivation. Sort and product conclusions have sort types.
Variable supplies the annotation's sort judgment, carried to its
extended context by weakening. Weakening carries the type-correctness
judgment from its subject premise; abstraction supplies the product's
sort judgment; conversion supplies the new type's sort judgment.
For application, apply the induction hypothesis to its function premise.
Its literal product type is not a sort, so it is well typed at some sort.
W2 applied to that product yields its domain and codomain sort premises;
the codomain sort is the product's sort, since sorts do not beta-convert.
The separately supplied substitution theorem, with the actual argument
premise, yields \(\Gamma\vdash C[Q/x]:s_2\). This proves the remaining
case. No subject-reduction hypothesis entered W3.

\paragraph{W4. Context conversion.}
Suppose \(\Gamma,x:B,\Theta\) is valid,
\(\Gamma\vdash B':s'\), and \(B=_\beta B'\).
Every judgment in the old context is derivable in
\(\Gamma,x:B',\Theta\), with its subject, type, and later raw
annotations unchanged. The new context is valid.
Use induction on the source derivation, split at the selected declaration.
Premises before \(x\) are unchanged, by the free-variable discipline
already established in the supplied substitution proof. At the selected
variable case derive \(x:B'\), then convert to \(x:B\): the original
validity certificate \(\Gamma\vdash B:s\) is available in the new
extension by weakening. At a later variable or weakening step transform
its annotation judgment first, making the new prefix valid, then apply
the same variable or weakening rule. Product and abstraction transform
the base and fresh extended premises and keep the same sort pair;
application transforms both premises; conversion transforms its subject
and result-type premises and keeps its raw equality. A sort case before
the selected declaration is unchanged. These cases also construct all
later annotation certificates, proving new-context validity. Reversing
the roles of \(B,B'\) proves the reverse transport when both are typed
in the prefix. This does not require subject reduction.

\paragraph{W5. Stratification and the complete source disambiguation.}
Choose disjoint variable classes \(V_D,V_L,V_P\). Domain expressions and
object expressions are the following auxiliary syntactic classes:
\[
 D ::= \mathrm{Prop}\mid a\ (a\in V_D)\mid D\to D,
\]
\[
 O ::= x\ (x\in V_L)\mid\lambda x:D.O\mid OO
       \mid\Pi x:D.O\mid\Pi z:O.O,
\]
where object binders belong to \(V_L\), proof binders to \(V_P\), and
the last object's codomain contains no occurrence of its proof binder.
The notation \(D\to E\) is its actual source product with a fresh,
unused binder; it is not an additional constructor. This broad object
grammar is not a typing claim: only source-typed members will be used.
In particular its product clauses need not be well typed without the
source formation premises.

Every \(D\) is beta normal. Every beta reduct of an \(O\) is again an
\(O\), including reduction in annotations. To prove the latter, first
prove structurally that substituting an object for an object variable
preserves \(O\): domain annotations contain no object variables;
fresh binders permit induction on bodies; a proof binder has no occurrence
in an object body. A root object beta redex thus has an object contractum,
and compatible subterm contraction preserves the grammar. The classes
\(D,O,\{\mathrm{Type}\},\{\mathrm{Type}'\}\) are pairwise disjoint:
the variable classes differ; products in \(D\) end recursively in a
domain atom, whereas products in \(O\) have an object body; abstractions
and applications occur only in \(O\); the remaining sorts are neither.
Consequently no members of different such classes are beta-equal, by
their reduction closure and raw confluence. Beta-equal domain expressions
are literally equal up to their irrelevant binder names.

For each valid source context, assign its variables incrementally to the
three classes by its annotation certificates. Every source derivation
then has exactly one of these four forms:
\begin{enumerate}
\item \(\Gamma\vdash\mathrm{Type}:\mathrm{Type}'\);
\item \(\Gamma\vdash D:\mathrm{Type}\), with \(D\) in the domain grammar;
\item \(\Gamma\vdash O:D\), with \(D\) in the domain grammar and
      \(\Gamma\vdash D:\mathrm{Type}\);
\item \(\Gamma\vdash p:F\), with \(F\) in the object grammar and
      \(\Gamma\vdash F:\mathrm{Prop}\).
\end{enumerate}
Here a declaration is respectively a domain declaration \(a:\mathrm{Type}\),
an object declaration \(x:D\), or a proof declaration \(z:F\).
The third case includes formulas themselves, at the domain
\(D=\mathrm{Prop}\). The fourth case imposes no new grammar on the
proof subject; its stated formula typing is carried as part of the invariant.

Prove context classification and this four-case invariant simultaneously
by induction on annotation certificates and typing derivations. The sort
axioms give cases two and one. In a variable rule, the annotation premise
at a sort can be only case one (annotation Type), case two (annotation
domain), or case three at domain Prop (annotation formula); the new
variable consequently has its prescribed domain/object/proof class.
Weakening assigns that class and retains the invariant for its subject.
For product formation with result sort Type both domain and codomain
are domain expressions. Its binder is an object variable and so cannot
occur in its domain-grammar codomain; this is a nondependent domain arrow.
With result sort Prop, the codomain is an object expression at domain
Prop. A Type-sorted annotation makes an object binder and gives
\(\Pi x:D.F\); a Prop-sorted annotation makes a proof binder, which
cannot occur in the object codomain, and gives \(\Pi z:F.G\).
Thus both admitted cases are in the object grammar at domain Prop.

For abstraction, its well-formed product premise has one of precisely
these product forms. A domain-arrow type requires a domain annotation
and object body, so the resulting lambda belongs to the object grammar.
A propositional-product type gives case four, and that product's
Prop typing is already its premise. For application with a domain-arrow
function type, the domain and codomain are rigid domain expressions;
function and argument are objects and the application is an object at
the unchanged, nondependent codomain. With a propositional function type,
W2 gives a Prop-sorted object codomain. An object argument substitutes an
object for an object variable, preserving that grammar. A proof argument
substitutes for a variable absent from the object codomain, leaving it
unchanged. The supplied substitution theorem gives the resulting formula's
Prop judgment, establishing case four. These cases exhaust applications.

For conversion, its well-typed target at a sort is either Type, a domain,
or an object formula, by the induction hypothesis. Its old type belongs
to one of the four separated type classes. Raw confluence therefore
forces the same class. A domain target and old domain are identical,
so the object invariant persists. Formula conversion retains case four
with the new target's explicit Prop judgment. Conversion to Type retains
case two; conversion from the exceptional Type' result is impossible,
since no well-typed target belongs to that class. This is the step that
rules out spurious syntax allowed by arbitrary untyped beta expansions;
neither a normalization theorem nor untyped expansion closure is used.

The induction proves more than free-variable absence: domain-expression
judgments reconstruct in the domain subcontext alone, and object judgments
reconstruct in its domain-and-object subcontext, in their original relative
order. For domains use their structural formation: a domain atom has its
axiom/projection, and a domain arrow uses Type/Type with a weakened codomain.
For objects use the same induction cases above. Object projection retains
its object declaration. Object abstraction and application retain just
their object premises. A universal-product clause retains the object
binder and its two premises. An implication-product clause first reconstructs
its two object formulas without the proof binder, then weakens the
codomain by that binder solely to apply the source Prop/Prop formation
rule. An object conversion has identical rigid domain types and can be
omitted. Every ambient proof weakening is omitted; domain/object weakening
is retained. This supplies actual smaller-context derivations rather
than appealing to a general unproved strengthening principle.

Now form \(\Gamma_D\) from all domain declarations, then
\(\Gamma_L\) from all object declarations, then \(\Gamma_P\) from all
proof declarations, keeping relative order within each. Domain declarations
have the constant annotation Type, so their reordered context is valid.
Every object annotation is a domain formed in \(\Gamma_D\); therefore
\(\Gamma_D,\Gamma_L\) is valid. Every proof annotation is a formula
reconstructed in the old domain/object subcontext; W1 carries it into
\(\Gamma_D,\Gamma_L\), so \(\Gamma_P\) can be appended. The permutation
is valid and has exactly the old declarations, hence W1 carries every
old judgment to it. All assertions of source Lemma~3.11 follow:
domains are simple arrows formed in \(\Gamma_D\); objects need no proof
declarations; a Prop/Prop product's body does not contain its proof binder.
In particular Type/Type dependence is impossible as well.

\paragraph{W6. Uniqueness of types.}
For the exact source system, \(\Gamma\vdash M:A\) and
\(\Gamma\vdash M:B\) imply \(A=_\beta B\).
Induct structurally on the common subject \(M\), using W2 to remove
terminal weakening/conversion in both derivations. Variable and sort
types are fixed by their declaration or single axiom. A product's two
codomain sort premises type the same smaller body in the same literal
extended context; the induction hypothesis makes these sorts beta-equal
and hence equal. Their product result sorts are therefore equal.
For an abstraction the annotation is the same syntactic domain in both
derivations. The body induction gives equality of the two body types;
compatible products give equality of the generated abstraction types.
For an application the function induction gives equality of its two
generated product types. W2's product injectivity gives beta equality of
their domains and bodies (with the same fresh binder). Substitution of
the common argument preserves beta equality, so the result types are
equal. Finally compose with the two generated outer type equalities.
These are every possible subject constructor. The proof needs only
generation, raw confluence, and the exact functional sort/product rules;
it assumes neither subject reduction nor checker uniqueness.

\paragraph{W7. Subject reduction, including annotation contraction.}
If \(\Gamma\vdash M:T\) and \(M\longrightarrow_\beta M'\), then
\(\Gamma\vdash M':T\). Induct structurally on \(M\), with the induction
statement quantified over every typing derivation and valid context for
that smaller subject. This avoids treating generated witnesses as if
they necessarily had smaller derivation heights. W2 yields a constructor
judgment at generated type \(T_0=_\beta T\).
If the subject reduces, W5 excludes its rigid domain and sort cases;
it is an object or proof, whose type \(T\) has a sort by W3 and W5.
Thus once the reduced constructor has type \(T_0\), conversion to
the original \(T\) is licensed. The following are the constructor cases.

For \(M=\Pi x:A.B\), a domain-annotation contraction gives
\(\Gamma\vdash A':s_1\) by the smaller-subject induction on \(A\).
W4 carries \(B:s_2\) to \(\Gamma,x:A'\), and the same allowed pair
rebuilds the product at \(s_2\). A body contraction uses the induction
on \(B\) under the old extension and rebuilds the same product.
For \(M=\lambda x:A.P\), W2 applied also to its well-formed product
gives \(A:s_1\), \(B:s_2\). A body contraction uses the induction
on \(P:B\). An annotation contraction uses the induction on
\(A:s_1\), then W4 to carry both \(P:B\) and \(B:s_2\) to the
new binder context. Product formation and abstraction give the new
lambda type \(\Pi x:A'.B\), beta-equal to the old generated type.
The original product type is well formed, so conversion restores it
before any outer conversion. This explicitly preserves typed annotations.

For \(M=PQ\), a contraction in either proper child uses its
smaller-subject induction with the generated product or domain type,
then the same application rule. At a root redex
\((\lambda x:A.P)Q\), generation of the function judgment yields
\(\Gamma,x:A\vdash P:B\), a well-formed product
\(\Pi x:A.B\), and its equality to the application's product
\(\Pi x:C.E\). Product injectivity gives \(A=_\beta C\),
\(B=_\beta E\). Its formation gives \(\Gamma\vdash A:s_1\),
so convert the actual argument \(Q:C\) to \(Q:A\). Substitution
gives \(\Gamma\vdash P[Q/x]:B[Q/x]\). The generated application
type \(E[Q/x]\) has a sort by the carried typing invariants in W5:
a domain is formed at Type, and an object formula is typed at Prop.
This also covers the literal domain Prop, for which W3's weaker
``type is a sort'' alternative alone would not supply the sort judgment.
Substitution preserves \(B=_\beta E\), so convert to
\(E[Q/x]\), then to \(T\). Sorts and variables have no contraction.
These cases cover beta at every source-compatible position.
Induction on a finite path yields multi-step subject reduction.

\paragraph{W8. A valid normalization bridge, not the printed inverse.}
Let the target be precisely source CC:
sorts \(\ast,\Box\), axiom \(\ast:\Box\), and all four product
pairs \((\ast,\ast),(\ast,\Box),(\Box,\ast),(\Box,\Box)\).
Define a homomorphism \(h\) on raw pseudo terms by
\(h(\mathrm{Prop})=h(\mathrm{Type})=\ast\),
\(h(\mathrm{Type}')=\Box\), fixing variables and the product,
abstraction, and application constructors. It commutes with substitution
by structural induction and preserves each designated beta contraction,
including contractions in annotations. It is not a PTS morphism:
the source axiom Prop:Type would homomorphically become the false
target axiom \(\ast:\ast\).

For a domain expression \(D\), define \(\sigma(D)=\ast\) when its
rightmost atom is a declared domain variable and \(\sigma(D)=\Box\)
when that atom is Prop. Equivalently, \(\sigma(a)=\ast\),
\(\sigma(\mathrm{Prop})=\Box\), and
\(\sigma(D\to E)=\sigma(E)\).
Map source contexts homomorphically on annotations; in particular
\(a:\mathrm{Type}\) becomes \(a:\ast\). Map their classified
judgments as follows:
\[
\begin{array}{rcl}
 \Gamma\vdash\mathrm{Type}:\mathrm{Type}'
   &\mapsto& h\Gamma\vdash\ast:\Box,\\
 \Gamma\vdash D:\mathrm{Type}
   &\mapsto& h\Gamma\vdash hD:\sigma(D),\\
 \Gamma\vdash O:D
   &\mapsto& h\Gamma\vdash hO:hD,\\
 \Gamma\vdash p:F
   &\mapsto& h\Gamma\vdash hp:hF.
\end{array}
\]
The result-sort override is part of this typed judgment map; it is not
an assumption that the homomorphic sort image types every image term.

Prove valid target contexts and these judgment translations simultaneously
by induction on the W5-classified source derivations. Both source sort
axioms become the target axiom \(\ast:\Box\). A domain declaration
becomes a type variable licensed by that axiom. An object declaration
is licensed by \(hD:\sigma(D)\); a proof declaration is licensed
by \(hF:\ast\). Variable and weakening consequently work at the
translated sorts. A source domain product becomes a target product with
domain sort \(\sigma(D)\), codomain sort \(\sigma(E)\), and result
sort \(\sigma(E)\). All four possibilities are admitted by exact CC.
A formula product has target codomain sort \(\ast\); its domain
sort is \(\sigma(D)\) for a data annotation, or \(\ast\) for a
formula annotation. The corresponding \((\ast,\ast)\) or
\((\Box,\ast)\) pair is admitted. Abstraction uses the translated
body and exactly this translated well-formed product. Application uses
the translated product and argument; substitution commutation identifies
the target result with the mapped source result.

In a source conversion of an object/proof type, the target type has its
translated sort judgment and the raw beta equality maps by \(h\), so
target conversion is legal. A source conversion in a domain judgment
retains the literal result Type by W5's separation, and the source domain
is unchanged; its already translated \(hD:\sigma(D)\) is retained
without an invalid attempted conversion to \(\ast\). The exceptional
top-result judgment has no conversion case, again by W5. These are all
rules, establishing transport on every actual Figure~2 derivation.

For example, with \(a:\mathrm{Type}\), the source domain
\((a\to\mathrm{Prop})\to a\) maps to
\((a\to\ast)\to a:\ast\). Its domain has target sort \(\Box\)
and its codomain has target sort \(\ast\), using exactly the CC pair
\((\Box,\ast,\ast)\). This is the mixed case missed by the printed
PRED-omega inverse described next.

It follows that strong normalization of this exact CC, if separately
supplied for its full compatible annotated beta reduction, entails strong
normalization of the exact source lambdaHOL. Indeed a hypothetical
infinite source path starts with a source-typed term (or source-typed
type); the translated starting subject is CC typed by the construction,
and every contraction maps to one contraction at the same position.
Thus it gives an infinite target contraction path, contrary to that
theorem. A source top sort appearing only as a type maps to \(\Box\)
and is already a beta-normal atom. Identifying two sort symbols erases
no lambda/application redex. This proves the transfer, not the target
CC normalization theorem itself.

\paragraph{W9. Exact CC normalization construction: source and standing.}
The following is an independently worded reconstruction of the direct
two-interpretation argument in Geuvers's December~1994 Computing Science
Report~94/50, sections~2--3, printed pp.~2--10. This is a separately
identified report manifestation, not a claim of byte identity with the
1995 LNCS publication cited by the 2001 chapter. The report's CC rules
are exactly the two-sort Figure~3 rules required by W8. Here the
candidate notion is the reduction-closed neutral-expansion version,
rather than the report's key-redex saturated-set version; every closure
property actually used is proved below. The argument uses ordinary sets
of finite syntax, powersets and finite iterated function sets; it needs
no inaccessible cardinal or set-theoretic model of an impredicative
universe as its own member. No CC normalization theorem is a premise.

\paragraph{W9a. CC structural prerequisites, with the untyped top case.}
For the exact CC rules, reapply W1's validity/projection/thinning,
W2's generation, W3's type-correctness, W4's context conversion, and
W6's structural uniqueness inductions. Their respective rule cases use
only the declared axiom, the actual admitted product pair, substitution,
and raw product injectivity; replace the two lambdaHOL axioms by
\(\ast:\Box\) and use the four exact CC pairs. In particular the
application type-correctness case still explicitly derives its result
type at the codomain's sort by substitution. In the uniqueness product
case the result is again the codomain sort, so equality of the two body
sorts gives equality of the result sorts. Nothing in these five inductions
uses the lambdaHOL-specific domain/object grammar of W5.

For precision, CC subject reduction is also supplied here rather than
assumed from W7, whose domain argument was lambdaHOL-specific.
Use structural induction on the subject, universally over its typings.
For product and lambda annotation contractions, use the smaller
annotation judgment, transport the body judgment by context conversion,
and rebuild the same admitted product or abstraction. Body contractions
use the smaller body judgment. In an application, proper-child contractions
use the generated function/argument premises. At a root lambda redex,
generation and product injectivity convert the actual argument to the
literal lambda annotation; substitution types the contractum at the
literal body instance; the generated application's result type is
explicitly well typed by the application case of type correctness, so
conversion restores it. These are the W7 construction steps, but here
the existence of that result-type sort is supplied directly by CC type
correctness, not by W5.

Outer conversion is licensed for every original type other than
\(\Box\): a nonsort has its sort by type correctness, and the only
other sort \(\ast\) has type \(\Box\) by the axiom and thinning.
For original type \(\Box\), no final conversion can have that target,
since no judgment types \(\Box\); generation proves this absence by
induction on the typing rules. Strip final weakenings. A variable of
literal type \(\Box\) would require a forbidden annotation certificate;
a lambda has a literal product type; an application of literal result
\(\Box\) contradicts its explicitly derived result-type sort judgment.
Thus the subject is \(\ast\), which does not reduce, or a product whose
formation result is literally \(\Box\). Product contractions rebuild
that same sort without an illicit conversion premise. This exhausts the
exception and establishes full compatible CC beta subject reduction.
Together with raw confluence it supplies typed common reducts; no
normalization was used. Raw confluence here is the main author's
parallel annotated-grammar proof with the sort alphabet instantiated
to \(\{\ast,\Box\}\), and substitution is the supplied Figure~3
whole-context theorem, not a new assumed CC normalization result.

If \(\Gamma\vdash A:s\), \(\Gamma\vdash B:t\), and \(A=_\beta B\),
with \(s,t\in\{\ast,\Box\}\), typed common reducts and uniqueness
give \(s=_\beta t\), hence \(s=t\). This is sort discrimination.
Also a typable subject \(\Gamma\vdash A:T\) cannot beta-convert to
\(\Box\): since that atom is normal, confluence would make it a
reduct, and subject reduction would type it. This excludes subjects,
not the normal atom \(\Box\) counted as well typed merely because
it occurs on the right of the axiom. These facts justify classification and
semantic conversion cases below independently of normalization.

\paragraph{W9b. CC classification and kind shapes.}
Partition variable names into constructor variables \(\alpha\) and
object variables \(x\), by the unique sort of each declaration's
annotation, renaming binders consistently. Write \(K,C,O\) for these
auxiliary syntax classes:
\[
\begin{aligned}
 K&::=\ast\mid\Pi x:C.K\mid\Pi\alpha:K.K,\\
 C&::=\alpha\mid\lambda x:C.C\mid\lambda\alpha:K.C
       \mid CC\mid CO\mid\Pi x:C.C\mid\Pi\alpha:K.C,\\
 O&::=x\mid\lambda x:C.O\mid\lambda\alpha:K.O\mid OO\mid OC.
\end{aligned}
\]
These grammars classify syntax, not arbitrary untyped grammar members
as typable. They are pairwise disjoint. For products, compare the body
class recursively; for lambdas likewise compare the body; for applications
compare the left child recursively; variables and the atom \(\ast\)
separate the initial cases. Matching substitution preserves all three
classes: a constructor substituted for a constructor variable, or an
object for an object variable, preserves \(K,C,O\). Prove this by
simultaneous structural induction: the selected variable is replaced
in its own class; application combines its child classes; fresh binder
renaming permits the annotation/body inductions in every binder case;
the atom is unchanged.

Every typed CC subject belongs to exactly the appropriate one of these
classes: \(M:\Box\) gives \(M\in K\); \(M:A\) with
\(A:\Box\) gives \(M\in C\); \(M:A\) with \(A:\ast\) gives
\(M\in O\). In particular an expression of type \(\ast\) is a
constructor, and a kind is an expression of type \(\Box\).
Prove the classifications simultaneously by induction on derivations,
carrying the literal class of the subject's type and the context tags.
The axiom supplies \(\ast\in K\). An annotation at \(\Box\) is a
kind, so its variable is a constructor; an annotation at \(\ast\)
is a constructor type, so its variable is an object. Weakening retains
these facts. A product with codomain sort \(\Box\) has kind body and
is in \(K\), with annotation in \(C\) or \(K\) according to its
domain sort. A product with codomain sort \(\ast\) has constructor
body and is in \(C\), with exactly the analogous two annotation cases.
For a lambda, its product-sort premise exposes the codomain sort. A
kind codomain requires a constructor body; a type codomain requires an
object body. This gives the two lambda-body classes in the grammar.

For an application, the function induction classifies its literal
product type. A kind product has kind body and constructor function;
a type product has constructor body and object function. The argument
induction gives constructor or object according to the domain's sort.
The matching-substitution lemma makes the resulting kind/type instance
remain in its class. The explicit substitution sort judgment supplies
its typing. Conversion preserves the sort of every typed old/new type
by W9a's sort discrimination. The untyped old top \(\Box\) cannot
convert to a well-typed target, again by W9a. Thus conversion retains
the induction class. This treats all primitive rule and sort cases,
including the top case, without an appeal to the report's unproved
classification citation. An expression \(\Box\) itself occurs only
as the exceptional type, never inside a typed subject.

Define a finite kind shape by
\(\kappa(\ast)=\mathbf b\),
\(\kappa(\Pi x:A.B)=\kappa(B)\) for an object binder, and
\(\kappa(\Pi\alpha:A.B)=\kappa(A)\to\kappa(B)\) for a constructor
binder. Kind syntax has no variable, lambda or application at its root;
constructor variables occur only inside its ignored type annotations.
Matching constructor/object substitution therefore leaves its shape
unchanged, by structural induction. A typed kind contraction also leaves
its shape unchanged: it cannot be a root redex; contractions in type
annotations are ignored; contractions in kind annotation/body use their
structural induction. The domain's type/kind tag cannot change under
that contraction by subject reduction and sort discrimination.
Beta-equal typed kinds consequently have identical shapes, by their
typed common reduct. These statements require no normal form for a kind.

\paragraph{W9c. Candidates and the actual beta-expansion lemma.}
Use finite untyped lambda terms with one inert constant \(d\), ordinary
compatible beta contraction, and \(\mathsf{SN}\) as the set having no
infinite such path. A candidate \(R\) is a subset of this one raw-term
set satisfying: (i) \(R\subseteq\mathsf{SN}\); (ii) it is closed
under contraction; (iii) a neutral term whose every immediate reduct is
in \(R\) belongs to \(R\). Here variables, \(d\), and applications
are neutral; only a lambda is nonneutral. A redex application is neutral
and is not exempted from (iii). Candidates exist: \(\mathsf{SN}\)
itself satisfies these clauses, since any infinite path from a neutral
term must start at an immediate reduct. Every candidate contains every
variable and \(d\), since they have no reducts.

For a strongly normalizing finite term \(P\), let \(\eta(P)\) be its
greatest reduction length. This is finite: there are finitely many redex
positions at each finite term; if arbitrarily long paths existed, at
each finite branching one could choose a child admitting arbitrarily
long continuations, constructing an infinite path. Reducts strictly
lower \(\eta\). Subterms of an SN term are SN, since any subterm path
lifts through its fixed context. An inert-headed application
\(dP_1\cdots P_n\) is SN when all \(P_i\) are SN: every step lowers
\(\sum_i\eta(P_i)\). A lambda is SN exactly when its body is SN.

The controlled expansion lemma is: if \(F,N\in\mathsf{SN}\) and
\(F[N/x]\in R\), then \((\lambda x.F)N\in R\).
Induct on \(\eta(F)+\eta(N)\). Its immediate root contraction is
the given member. A contraction in \(F\) gives \(F'\); substitution
preserves that contraction, so \(F'[N/x]\in R\) by reduction closure,
and the smaller induction gives \((\lambda x.F')N\in R\).
A contraction \(N\to N'\) gives
\(F[N/x]\to^*F[N'/x]\) by reducing its finitely many substituted
copies, possibly zero, so reduction closure and the smaller induction
give \((\lambda x.F)N'\in R\). These are all immediate reducts;
neutral expansion proves the desired membership. Substitution preservation
and the finite-copy statement follow structurally on untyped syntax,
freshening binders, by the same capture-avoiding substitution identity
as in the main author's substitution proof. This lemma does not assert
arbitrary beta expansion preserves SN.

Nonempty intersections of candidates are candidates, clause by clause.
The arrow
\(R\Rightarrow S=\{F:\forall N\in R,\ FN\in S\}\)
is also a candidate. First choose a fresh variable in \(R\): if
\(F\) were not SN, its path would lift into \(Fx\in S\), a
contradiction. Reduction closure follows by lifting \(F\to F'\)
into each application and using closure of \(S\). For neutral
expansion, suppose every immediate reduct of neutral \(F\) belongs
to the arrow. For \(N\in R\), induct on \(\eta(N)\).
Every immediate reduct of \(FN\) is either \(F'N\), in \(S\)
by the assumption, or \(FN'\), in \(S\) by closure of \(R\) and
the smaller induction. There is no root contraction because \(F\)
is not a lambda. Thus \(FN\in S\) by neutral expansion. This proves
the last candidate clause and finishes arrow closure.

Arrow introduction is now justified: if \(F[N/x]\in S\) for every
\(N\in R\), then \(\lambda x.F\in R\Rightarrow S\).
Take \(N=x\), after the usual fresh choice, to get \(F\in S\)
and hence its SN premise; all other \(N\) are SN by membership in
\(R\). Controlled expansion gives the required applications.
Let \(\mathrm K=\lambda p.\lambda q.p\). Applying this introduction
twice shows \(\mathrm K\in R\Rightarrow(\mathsf{SN}\Rightarrow R)\)
for every candidate \(R\). Hence
\(F\in R,\ A\in\mathsf{SN}\) implies \(\mathrm KFA\in R\).
This is the annotation-retention lemma used below, with both premises
proved rather than hidden inside an expansion assumption.

\paragraph{W9d. Finite semantic domains and canonical inhabitants.}
Let \(\mathcal C\) be the set of all candidates, a subset of the
powerset of the fixed raw-term set. Interpret a finite kind shape by
\(V(\mathbf b)=\mathcal C\) and
\(V(k\to\ell)=V(\ell)^{V(k)}\), the ordinary function set.
Write \(V(A)=V(\kappa(A))\) for a source kind. Each is a set by
finite recursion, and the collection over all finite shapes is a set
by replacement over the countable finite-shape syntax. No universe
closure axiom is used. Each is nonempty with explicit canonical element:
\(c_{\mathbf b}=\mathsf{SN}\), and
\(c_{k\to\ell}\) is the constant function of value \(c_\ell\).
Kind-shape invariance gives \(V(A[P/\alpha])=V(A[t/x])=V(A)\)
and \(V(A)=V(B)\) for beta-equal typed kinds. These equalities do not
say that all kind candidates are equal.

A semantic constructor valuation \(\xi\) assigns to every declared
\(\alpha:A\) an element of \(V(A)\). These domains do not depend
on object assignments or on \(\xi\). Define a partial structural
interpretation for constructors and kinds, and an auxiliary interpretation
of the exceptional top, by
\[
\begin{aligned}
 [\ast]_\xi&=[\Box]_\xi=\mathsf{SN},&
 [\alpha]_\xi&=\xi(\alpha),\\
 [PQ]_\xi&=[P]_\xi([Q]_\xi)&&\text{for a constructor argument }Q,\\
 [Pt]_\xi&=[P]_\xi&&\text{for an object argument }t,\\
 [\lambda\alpha:A.P]_\xi&=(a\mapsto[P]_{\xi[\alpha:=a]})
      &&(a\in V(A)),\\
 [\lambda x:A.P]_\xi&=[P]_\xi&&\text{for an object binder},\\
 [\Pi x:A.B]_\xi&=[A]_\xi\Rightarrow[B]_\xi
      &&\text{for an object binder},\\
 [\Pi\alpha:A.B]_\xi&=[A]_\xi\Rightarrow
           \bigcap_{a\in V(A)}[B]_{\xi[\alpha:=a]}
      &&\text{for a constructor binder}.
\end{aligned}
\]
The product codomain \(B\) is a type or kind, so its interpretation is
a candidate, not an arbitrary semantic function. The constructor-lambda
body \(P\) is a constructor, whose interpretation may be a function.
These distinctions are supplied by W9b. An object does not itself have
a set interpretation; it is only ignored when occurring as an argument
of a constructor. Thus \([\ast]\) is a candidate, whereas its kind
domain \(V(\ast)\) is the set of candidates; they are not identified.

\paragraph{W9e. Definedness, membership, and semantic substitution.}
For every actual typed constructor \(P:A\) with \(A:\Box\), and
every semantic valuation of its context, this interpretation is defined
and \([P]_\xi\in V(A)\). For every kind \(A:\Box\), its
interpretation is a candidate. In particular every type \(T:\ast\)
has a candidate interpretation. Prove these assertions simultaneously
by structural induction on the classified constructor/kind syntax,
quantified over its source typing derivations; W9a generation supplies
the typing cases. A constructor variable has the prescribed valuation
member. A kind atom gives the candidate SN. A constructor lambda over
a constructor binder gives a genuine function from \(V(A)\) to
\(V(B)\), since every extended semantic valuation is valid and the
body induction supplies its value there; these are exactly the target
kind shape's function sets. A constructor lambda over an object binder
retains the body value; its target kind domain also ignores that binder.

In a constructor application to a constructor, the function induction
gives a member of \(V(A)\to V(B)\), and the argument induction
gives its input in \(V(A)\), so the semantic application is defined
and lies in \(V(B)\). The actual result kind has this same domain by
kind-shape substitution. With an object argument, both the semantic
value and the result's kind domain are unchanged. In a product with
object binder, the two proper-subterm inductions give candidate domain
and codomain, so arrow closure gives a candidate. With constructor
binder, its proper kind-domain induction gives a candidate, and every
extended valuation gives a candidate codomain. The indexing set \(V(A)\)
is nonempty by W9d, so intersection and arrow closure give a candidate.
For a product at sort \(\ast\), this candidate is exactly membership
in \(V(\ast)=\mathcal C\). For a product at sort \(\Box\), it
is the asserted kind-candidate property. Conversion of a constructor's
kind type retains the semantic-domain membership by kind-shape equality
from W9b; conversion does not change the recursively evaluated subject.
A target literal \(\Box\) has no conversion case, as in W9a.
Weakening or a changed ambient valid context changes no used assignment.
These cases exhaust the syntax and its typing rules. This argument
establishes the evaluator's definedness; it does not assume it when
applying a purported undefined function.

The exact substitution equalities are
\[
 [B[P/\alpha]]_\xi=[B]_{\xi[\alpha:=[P]_\xi]},\qquad
 [B[t/x]]_\xi=[B]_\xi,
\]
for typed matching substitutions and a constructor or kind \(B\).
Prove them by simultaneous structural induction, with fresh binders.
A constructor variable gives its selected replacement or unchanged
valuation; the atom is unchanged. Constructor application uses the two
constructor substitutions and ordinary function application. An object
argument stays an object under matching substitution and remains ignored.
At a constructor lambda, the bound variable is fresh, its value ranges
over the unchanged \(V(A)\) by W9d, and the body induction gives the
same function pointwise. At an object lambda, both sides ignore the
binder and the body induction applies. At a product the annotation and
body inductions identify the arrow factors; for a constructor binder the
indexing domain is unchanged and the equality holds at every index, so
the intersections agree. Object substitution cannot introduce a new
constructor argument or alter a variable tag; matching-substitution
classification supplies this fact. These are all clauses.

The interpretation is invariant along every typed beta contraction.
At a constructor-binder root, its semantic function application and the
first substitution equality give the contractum. At an object-binder
root, both abstraction and argument were ignored, and the second equality
gives the contractum. A kind has no root redex. Compatible steps in
constructor children use these same identities and structural induction;
object-only children are ignored. A step in a binder annotation keeps
its type/kind tag by W9a, keeps the indexing domain by W9d when it is a
kind, and keeps its candidate factor by this smaller-subterm induction.
Consequently every binder clause is preserved, including product factors
and constructor-lambda function domains. Typed common reducts from W9a
then show \([A]_\xi=[B]_\xi\) for beta-equal well-typed types or
kinds. This is why untyped conversion is sound for this evaluator; no
normalization or untyped expansion closure of classified syntax is used.

\paragraph{W9f. The annotation-retaining untyped interpretation.}
Define \(E\) structurally on typed subjects, with \(d\) inert and
\(\mathrm K=\lambda p.\lambda q.p\):
\[
\begin{aligned}
 E(\ast)&=d,& E(v)&=v,& E(MN)&=E(M)E(N),\\
 E(\lambda v:A.M)&=\mathrm K(\lambda v.E(M))E(A),&&&\\
 E(\Pi v:A.B)&=dE(A)(\lambda v.E(B)).&&&
\end{aligned}
\]
Top \(\Box\) need not be interpreted as a subject, since W9b shows
it never occurs in one; for an isolated type it is already normal.
All binders are freshened for substituted images. Structural induction
proves \(E(M[N/v])=E(M)[E(N)/v]\): the variable/atom/application cases
are immediate; a binder uses the annotation/body inductions and the same
fresh name. The annotation is outside its binder in both source and
target, which preserves the exact source binding scope.

A designated root contraction is simulated by three target contractions:
\[
 \mathrm K(\lambda v.E(M))E(A)E(N)
 \longrightarrow^3 E(M)[E(N)/v]=E(M[N/v]).
\]
For a designated contraction in an annotation or body, its positive
induction path lifts through the actual corresponding target context.
For a contraction in an application child it likewise lifts through
application. The product annotation and body have both been retained
in \(dE(A)(\lambda v.E(B))\); the lambda annotation remains present
as an argument of K until its administrative contraction. Thus every
source compatible contraction, at every position, has a positive finite
target path. We do not erase an annotation first and then claim its
reductions survive. Concatenating the positive paths from an infinite
source path would give an infinite target path. Hence \(E(M)\in
\mathsf{SN}\) implies full compatible source SN, including annotations.

For a finite raw substitution \(\rho\) on the context variables, let
\(E_\rho(M)=E(M)\rho\), substituting once into the original syntax.
This definition and the preceding substitution identity give
\(E_\rho(M[N/v])=E_{\rho[v:=E_\rho(N)]}(M)\) when the bound name
is fresh for the other images. In particular the image of a binder
under \(\rho\) uses its body evaluated with its own variable left
as that variable, not recursively substituted through the inserted images.

\paragraph{W9g. Fundamental lemma: every sort/product/binder case.}
Call \((\xi,\rho)\) a realized valuation of \(\Gamma\) when
\(\xi\) is its semantic constructor valuation and, for every
declaration \(v:A\), \(\rho(v)\in[A]_\xi\). Every annotation is
a type or kind, so W9e supplies this candidate. For every actual typing
\(\Gamma\vdash M:T\) and every realized valuation,
\[
                         E_\rho(M)\in[T]_\xi.
\]
Prove this by structural induction on the subject, universally over its
typings, using generation and semantic conversion invariance. This
induction calls only proper annotation/body/application children, not
possibly larger synthesized types. A generated outer conversion replaces
\([T_0]\) by the identical \([T]\); the top case has no conversion.
Generation removes final weakenings, or restricts the valuation to the
used prefix. The constructor cases are as follows.

For the sort \(\ast:\Box\), the image is \(d\in\mathsf{SN}
=[\Box]\). A variable's image is its realized-valuation member of
the declared type candidate; conversion invariance gives its final type.
For a product \(\Pi v:A.B:s_2\), the annotation induction yields
\(E_\rho(A)\in[s_1]=\mathsf{SN}\), for either \(s_1\).
Extend the valuation by \(\rho(v)=v\). A fresh variable belongs to
every candidate. If this is a constructor binder, also choose the explicit
canonical \(c_{\kappa(A)}\in V(A)\) for \(\xi(v)\); if it is an
object binder no semantic assignment is added. Hence the body induction
gives \(E_{\rho[v:=v]}(B)\in[s_2]=\mathsf{SN}\), for either
\(s_2\). Its lambda is SN. The inert-headed application
\(dE_\rho(A)(\lambda v.E_{\rho[v:=v]}(B))\) is therefore SN,
which is membership in the product's target sort interpretation.
This covers all four product triples, including \((\Box,\Box,\Box)\),
not just the type/object product cases printed in the report's abbreviated
soundness argument.

For an abstraction over an object \(x:A\), the generated body is
\(P:B\) and its target product interpretation is
\([A]_\xi\Rightarrow[B]_\xi\). For every \(p\in[A]_\xi\),
the extended realized valuation gives by induction
\(E_{\rho[x:=p]}(P)\in[B]_\xi\); object assignments do not change
the semantic candidate of the codomain. By the E substitution identity,
these are exactly the instances of the fresh-body term
\(F=E_{\rho[x:=x]}(P)\). Arrow introduction from W9c gives
\(\lambda x.F\) in the required arrow. The annotation induction
gives \(E_\rho(A)\in\mathsf{SN}\); K retention therefore puts
the whole image of the annotated source abstraction in that arrow.
This holds whether \(P\) is a constructor or an object.

For an abstraction over a constructor \(\alpha:A\), the target
candidate is \([A]_\xi\Rightarrow J\), where
\(J=\bigcap_{a\in V(A)}[B]_{\xi[\alpha:=a]}\).
For every \(p\in[A]_\xi\) and every \(a\in V(A)\), extend both
valuations by \(\rho(\alpha)=p\), \(\xi(\alpha)=a\).
The body induction gives membership in the corresponding codomain
candidate. Its raw image depends only on \(\rho\), not on \(\xi\),
so the same term belongs to the entire intersection \(J\). It follows
by arrow introduction and annotation K retention that the whole lambda
image belongs to \([A]_\xi\Rightarrow J\). This argument also
covers both object and constructor bodies and supplies the SN body
premise using the fresh-variable instance; it never assumes the source
lambda was already normalizing.

For application with object argument \(N:A\), the generated function
candidate is \([A]_\xi\Rightarrow[B]_\xi\). Function and argument
inductions therefore give \(E_\rho(M)E_\rho(N)\in[B]_\xi\).
Semantic object-substitution invariance identifies this with
\([B[N/x]]_\xi\), the actual result type candidate.
For application with constructor argument \(N:A\), the function
candidate is \([A]_\xi\Rightarrow J\) with the preceding
intersection. Application of the two induction members gives membership
in \(J\). Independently W9e supplies \([N]_\xi\in V(A)\), so choose
that very index. The resulting membership is in
\([B]_{\xi[\alpha:=[N]_\xi]}=[B[N/\alpha]]_\xi\) by semantic
substitution. These two application cases cover both object and
constructor functions; no semantic value is assigned to an object
argument. Products, abstractions, applications, variable, and sort are
every subject constructor, so the fundamental lemma is established.

\paragraph{W9h. Open-context realization and full strong normalization.}
Every valid finite CC context has a realized valuation with raw
\(\rho(v)=v\): assign to each constructor declaration
\(\alpha:A\) the explicit canonical member
\(c_{\kappa(A)}\in V(A)\). Kind-domain shapes do not depend on any
prior semantic value, so these assignments are valid simultaneously.
W9e says every annotation's interpretation is a candidate; its variable
belongs to that candidate by W9c. Thus identity \(\rho\) realizes
all object and constructor declarations, including uninhabited source
types: a raw open variable is a candidate member, not a fabricated
closed proof of that type. No closure or inhabited-context assumption
has been added.

For \(\Gamma\vdash M:T\), the fundamental lemma now gives
\(E(M)\in[T]_\xi\subseteq\mathsf{SN}\). W9f's positive path
simulation implies that \(M\) has no infinite full compatible beta
path. If a well-typed pseudo term occurs only on the right of a judgment,
type correctness gives a sort judgment for it, unless it is a sort;
\(\ast\) has its axiom judgment and \(\Box\) is a normal atom.
Hence the same conclusion covers both subjects and types in the
source's definition of well typed. The proof treats arbitrary well-formed
open contexts, all four product triples, reductions beneath binders,
and contractions inside annotations/types. It assumes neither CC SN,
closed-term normalization, weak-head termination, checker completeness,
nor an impredicative set universe. Combining this constructed CC result
with W8 proves strong normalization for the exact beta-only source
lambdaHOL, leaving definitions, delta, recursors, and later inductive
extensions outside the conclusion.

\paragraph{W10. Failure of the source's printed inverse PRED-omega map.}
In the source context \(a:\mathrm{Type}\), Type/Type forms
\(a\to\mathrm{Prop}:\mathrm{Type}\) and then
\((a\to\mathrm{Prop})\to a:\mathrm{Type}\).
The printed judgment-dependent inverse maps \(a:\mathrm{Type}\)
to \(a:\mathrm{Set}\), with \(a\to\mathrm{Prop}:\mathrm{Type}\)
in PRED-omega. Its final arrow ending in \(a\) is assigned sort Set,
but its product formation would require
\((\mathrm{Type},\mathrm{Set},\mathrm{Set})\), which is absent from
the exact six source PRED-omega triples. Generation proves rejection:
the two rigid argument/codomain sorts cannot be changed by conversion.
Thus the displayed inverse is not derivation preserving on this legitimate
lambdaHOL domain. This countercase does not affect the separately proved
forward PRED-omega-to-lambdaHOL morphism, nor does it refute lambdaHOL
normalization. W8 supplies a different, explicitly editorial CC transport
that handles this omitted case without changing either source calculus.

\paragraph{W11. The live P7 interface and separate P12-draft/W9 bridge.}
This is a separate construction for the lead's later reconciliation, not
a change to the live P7 theorem. At the inspection checkpoint the live
route is \texttt{handbook/main.tex:35}, then
\texttt{handbook/parts/p07.tex:119}, then
\texttt{handbook/parts/p07/lambda-calculi-and-type-theories.tex:465}.
The result labelled \texttt{thm:p07-coc-strong-normalization} is
\texttt{cited=CoquandHuet1988} and asserts normalization of the exact
calculus presented by Coquand and Huet, excluding later inductive
extensions. Its live comparison-table consumer is at that chapter's
lines~567--569; the separate live proof-complexity consumer is
\texttt{handbook/parts/p07/proof-complexity.tex:300--301}.
The P7 annex driver selects no CoC normalization annex.

That P7 chapter selects no CoC pseudo-term grammar, sort specification,
context formation, inference sheet, conversion generators, or compatible
reduction relation. Its nearby \texttt{sheet:mltt-intensional} is not a
CoC sheet. Nor does its earlier unannotated untyped-lambda beta definition
declare a reduction relation for annotated CoC types. Therefore the
legacy citation is not an explicitly presented calculus to which one
can currently apply a checked bijective notation map. The original
1988 article's actual rule and equality extent has not been inspected
in this block. In particular no claim that it has, or lacks, eta, delta,
universe constraints, or a different context convention is inferred
from its title or from the 1994 report's attribution.

\paragraph{W11a. The genuine formal-P12-draft/W9 identity bridge.}
The current source-author P12 draft does declare an exact calculus:
\texttt{def:p12-har-pts-specifications}, with
\texttt{sheet:p12-har-pts}, specifies sorts \(\ast,\Box\), axiom
\(\ast:\Box\), and the four triples
\((\ast,\ast,\ast),(\ast,\Box,\Box),
(\Box,\ast,\ast),(\Box,\Box,\Box)\).
Its \texttt{thm:p12-har-cc-normalization} expressly selects full
compatible annotated beta reduction. At this checkpoint the P12 body and
annex drivers do not input that source-author presentation pair; it is
therefore not yet a reader-routed P12 result. This exact formal draft calculus
and W9's CC have a bijective bridge \(r\): fix both sorts and every
variable, preserve application, and send
\(\Pi x:A.B\) to \(\Pi x:rA.rB\) and
\(\lambda x:A.M\) to \(\lambda x:rA.rM\).
Send a context \(x_1:A_1,\ldots,x_n:A_n\) to the ordered context
\(x_1:rA_1,\ldots,x_n:rA_n\). The inverse uses the same clauses.
Structural induction proves both composites are identity, including
annotations, binder scopes and every context position. Alpha-equivalent
representatives map to alpha-equivalent representatives since variable
renaming is unchanged. The same structural induction gives
\(r(M[N/x])=rM[rN/x]\): a selected variable gives the replacement;
other variables and sorts are unchanged; application uses its two
inductions; a fresh lambda/product binder uses annotation and body
inductions. This supplies actual capture-avoiding substitution transport.

The root beta generator is preserved and reflected at exactly one step:
the image of \((\lambda x:A.M)N\) contracts to
\(rM[rN/x]=r(M[N/x])\), and every image root contraction has this
unique inverse source redex. A compatible step at any proper position
maps to a step at the same position: application children, lambda
annotation/body, and product annotation/body have been preserved.
Applying the inverse proves reflection in all these cases. Thus finite
beta paths, their reflexive/symmetric/transitive equality closure, and
infinite beta paths correspond in both directions. No redex is erased
and no eta/delta/iota generator is silently introduced.

Prove typing transport in each direction by induction on derivations.
The single sort axiom is unchanged. Variable formation uses its mapped
annotation-sort judgment and the unchanged fresh name. Weakening uses
its mapped annotation and subject judgments, keeping declaration order.
Product formation uses the two mapped sort judgments and precisely the
same one of the four listed triples. Abstraction uses the mapped body
judgment and mapped well-formed product. Application uses the mapped
function/argument judgments, with substitution commutation identifying
its result type. Conversion uses the mapped subject and well-sorted
target premises and the proved beta-equality correspondence. These are
every Figure~3 rule. The inverse induction has the identical seven
cases, so this is equivalence of judgments, not just a sound embedding.
Context validity corresponds both by these typed judgments and by the
equivalent prefix annotation certificates established in W9a; fresh
variables and every dependent annotation remain in their original order.
Hence the source's subject-or-type definition of well typed corresponds
as well. Finally an infinite normalization counterexample on either side
maps to an infinite one on the other. W9 therefore proves the formal
draft's exact beta-CC normalization claim, without a closed-context
restriction, and conversely both normalization claims have identical scope.

\paragraph{W11b. The precise legacy reconciliation boundary.}
Once the source-author P12 pair is actually routed, W11a allows the lead
to connect a P7 consumer explicitly scoped to that same two-sort beta PTS
to the new proved P12 result. It does not justify
replacing the legacy \texttt{cited=CoquandHuet1988} standing while keeping
the words ``exact calculus presented by Coquand and Huet'' unexamined.
The original source facet is the 1988 article, \emph{Information and
Computation}~76(2--3), pp.~95--120, DOI
\texttt{10.1016/0890-5401(88)90005-3}, identified by the canonical
bibliography entry. A bibliographic record is not its mathematical body.
The separate proof facet is the held Geuvers December~1994 report and
the editorial candidate reconstruction W9. The original facet remains
unresolved until its actual syntax, context discipline, rules, conversion
and full normalization relation are supplied, and a complete bridge to
the specified beta PTS is proved if appropriate. If that source uses
additional equality/reduction clauses, a matching preservation and
normalization argument for those clauses is still required; W11a has
no such clause in its domain.

The neighboring selected MLTT cannot close this interface. Its sheet
has separate context/type/term/equality judgments, Pi/Sigma/identity
constructors, unannotated term lambdas, and the cumulative Russell
axiom family \(\mathcal U_i:\mathcal U_{i+1}\). These are not W9's
two sorts, one axiom, and sole product/lambda/application grammar;
in particular an infinite primitive universe-sort family has no
sort-bijective renaming onto W9's two-element sort set. The selected
sheet itself assumes no additional Pi-eta law, whereas P7's separately
cited Abel--Coquand--Dybjer comparison concerns typed beta-eta equality
in its own subcalculus. Neither that comparison nor the CIC pointer
defines the cited 1988 CoC. No reader, P7, shared, or canonical edit is
made by this construction; the lead owns the editorial reconciliation.

\paragraph{W12. Qualified original 1988 body and the corrective full-system bridge.}
The original journal-layout Coquand--Huet article is now held, with its
printed \emph{Information and Computation}~76, pp.~95--120, 1988 heading,
1988 Academic Press copyright, and PII
\texttt{0890-5401(88)90005-3}. It has 26 consecutive printed pages.
The read scope is the base syntax and rules in sections~1--3, pp.~95--105,
the stripping/normalization standing in sections~4--5, pp.~105--110,
and the distinctions from variants/extensions in sections~6--8,
pp.~111--117. This is not a substituted 1985 or 1986 report.
The formal qualification record is
\texttt{spec/research/p07-coquand1988-source-qualification.md}.

The source has one written constant \(\ast\), one-based de Bruijn
indices, application, annotated lambda, and product \([x:A]B\).
Its syntactic logical contexts \(\mathcal A_o^n\) are \(\ast\)
or products whose codomain is another logical context. Its objects
\(\mathcal A_i^n\) are indices, applications of objects, lambdas with
object bodies, and products with object bodies; either class may be a
binder annotation. An environment is a closed logical context ending
in \(\ast\), viewed as an ordered list of declarations. The distinct
judgments are \(\Gamma\vdash\Delta\) for logical-context validity,
\(\Gamma\vdash M:P\) for object typing, and typed conversion
\(\Gamma\vdash M\cong N\). The source \(\ast\) has no type.
Thus no bijective sort renaming can turn this surface into
\(\ast:\Box\). Validity must instead become typing at the target's
auxiliary top \(\Box\). The source full variant is exactly the
starred-rule replacement definition on p.~102, not its restricted variant.

\paragraph{W12a. Literal-source corrections that must not be hidden.}
The p.~98 relocation is described as moving only free references up
one level, but its printed \(M^+=\xi_0(M)\), binder-body
\(\xi_{i+1}\), and variable condition \(k<i\) do not implement
that description with the expressly one-based indices. For the closed
legal object \(L=(\lambda x:\ast)1\), literal \(\xi_0\) gives
\((\lambda x:\ast)2\), introducing an external reference where the
source had only a bound one. Hence a literal use of this recurrence
does not commute with the intended binding interpretation. The
free-reference-only relocation used in the bridge below is the
explicit correction \(k\le i\) in that variable condition; with
\(i=0\) at the root, indices at most \(i\) are bound and unchanged,
and greater indices increase. An equivalent repair starts the cutoff
at one and retains the printed strict comparison. No correction is
made to the held primary PDF.

Also the restricted starred rule on p.~101 literally concludes
\(\Gamma\vdash M\cong\Delta\) from \(\Gamma\vdash M:\Delta\).
In \(\Gamma=[A:\ast]\), projection gives \(A:\ast\), and that
printed rule gives \(A\cong\ast\). Its named images are distinct
beta-normal terms, so this is not a beta-sound equality rule. The
same example contradicts the source's Lemma~6 on p.~105: its logical
rank assigns zero to object \(A\) and one to \(\ast\), although
that literal equality rule relates them.
p.~102 full replacement instead prints ordinary reflexivity
\(M:N\) implying \(M\cong M\). The bridge uses that actual full
replacement and never imports the bad restricted rule. Thus the
uncorrected restricted system is not silently called a subsystem of
the full system. The conclusions below concern the full rule surface
with the explicitly corrected one-based relocation, or equivalently
its stated concrete binding semantics. They are not an equivalence
claim for the literal misprinted indexed surface.

\paragraph{W12b. Indexed-to-named syntax, environments, and substitutions.}
For an ordered external variable vector \(\vec z=(z_1,\ldots,z_n)\)
from oldest to newest, define \(T_{\vec z}(\ast)=\ast\) and
\(T_{\vec z}(k)=z_{n+1-k}\). Map application homomorphically.
Map a source product to \(\Pi z:T_{\vec z}(A).
T_{\vec z,z}(B)\), and a lambda to the analogous annotated lambda,
choosing \(z\) fresh for the external vector. An environment
\(\Gamma=[x_n:A_n]\cdots[x_1:A_1]\ast\) maps to the ordered
declarations whose annotations are translated in their respective
older prefixes, rather than to a single type term. A logical context
appearing as a type \(\Delta\) still maps by \(T\) to a product
ending in \(\ast\). These are distinct roles of the source's
same product syntax and are not conflated.

The named encoding has a unique indexed inverse on its image, modulo
alpha-renaming: each variable names its position in the external vector
or binder stack; products/lambdas extend that stack only for the body,
and application recursively decodes both children. Induction on syntax
proves both composites are identity on the scoped source grammar and
its named image. It is not claimed to be onto arbitrary raw CC terms,
which additionally contain \(\Box\) and many ill-typed forms.

Corrected relocation commutes with this encoding. More generally, raise
the external indices by \(k\) with the correct cutoff under every
internal binder. A variable then still designates the same older named
variable; an internal bound reference still designates its binder.
The star/application cases are immediate; annotation/body induction
proves the binder cases. In particular the source lookup
\(\Gamma/l=\operatorname{add}(l,A_l)\) translates to the
original named annotation of the corresponding declaration, not to
an annotation containing that declaration's own fresh name.

For source \(Q\) scoped over \(\Gamma[x:A]\) and \(N\) scoped
over \(\Gamma\), corrected substitution/removal of index one gives
\[
 T_{\vec z}(Q[N/1])
   =T_{\vec z,z}(Q)[T_{\vec z}(N)/z].
\]
For index one this gives the argument; a larger external index decreases
by one when the removed binder disappears, identifying the same external
name on both sides. Under an inner binder, its local references stay
fixed, the substituted argument is correctly relocated, and a fresh
named binder prevents capture. Induction on annotation and body proves
both product and lambda cases, and application combines the child
identities. The same structural proof gives simultaneous substitutions
in all later environment annotations. A source beta contraction thus
maps to one designated named beta contraction at the same position,
including a contraction in a binder annotation or product body.

\paragraph{W12c. Joint full-source validity, typing, and equality transport.}
Write \(\widehat\Gamma\) for the named environment just defined.
Simultaneously induct on original full-system derivations to establish:
\(\Gamma\vdash\Delta\) gives
\(\widehat\Gamma\vdash_{\rm CC}T\Delta:\Box\);
\(\Gamma\vdash M:P\) gives
\(\widehat\Gamma\vdash_{\rm CC}TM:TP\);
and \(\Gamma\vdash M\cong N\) gives raw
\(TM=_\beta TN\), together with the following endpoint transport:
for every target type \(U\), either endpoint's CC judgment at \(U\)
implies the other endpoint's judgment at that same \(U\).
This stronger equality induction supplies the target sort premise needed
for source conversion; it does not assume the source's Lemmas~3--4,
whose full proofs the article refers elsewhere.

The primitive validity axiom maps to \(\ast:\Box\) in the empty
target environment. Extension by a source logical context \(\Delta\)
uses its mapped \(\Box\) judgment as the annotation certificate;
extension by an object \(M:\ast\) uses its mapped \(\ast\)
judgment. Both make the extended target environment valid, and the
weakened axiom types \(\ast\) there. A validity product rule has
premise \(\Gamma[x:A]\vdash\Delta\); its induction gives
\(T\Delta:\Box\) in the extended target context. CC validity
extracts the annotation certificate \(TA:s_1\); the triple
\((s_1,\Box,\Box)\) forms its image product. An object product rule
has body premise \(M:\ast\); similarly use
\((s_1,\ast,\ast)\). Both possible domain sorts are admitted.

For source projection, its validity premise makes
\(\widehat\Gamma\) valid. CC projection types the corresponding
named variable at its older-prefix annotation, and W12b's lookup
identity identifies this with \(T(\Gamma/l)\). A source abstraction
uses its body induction in the target extension. Its annotation sort
is obtained from that extension's CC validity. Its body type \(TP\)
has a sort by CC type correctness; if it is the literal sort \(\ast\),
use \(\ast:\Box\). It cannot be \(\Box\), which has no source
preimage. The requisite product is therefore well formed by one of
the four triples, and CC abstraction applies. Application uses the
two mapped premises and W12b's substitution identity.

Full equality reflexivity (for typed objects or valid logical contexts)
and symmetry/transitivity map to raw beta equality and preserve endpoint
typing immediately or by composing the endpoint inductions. For product
congruence, suppose an endpoint product is typed at \(U\). CC generation
exposes its annotation sort and body sort. Annotation equality induction
transports the first sort judgment to the second annotation; body
equality induction transports the body sort in the old binder context.
CC context conversion moves it to the new, beta-equal annotation.
Rebuild the product at its original result sort, then restore \(U\)
by conversion when required. If \(U=\Box\), the generated product
result is literally \(\Box\), so no illegal top-sort conversion is
used. For the reverse direction, first transport the new domain sort
to the old annotation, then use context conversion to carry the
second body judgment into the old binder context, where the body
equality induction is stated. Apply its reverse implication there
and rebuild the first product. This preserves the required body
context rather than assuming the equality was supplied in both contexts.

For lambda congruence, generation of a typed endpoint gives its literal
body judgment and a sorted product type. Annotation equality transports
its domain sort; body equality transports its body at the same generated
body type. Context conversion transports that body judgment and its
body-type sort to the changed annotation. Abstraction produces the new
product type, beta-equal to the old one; convert back to \(U\).
No lambda has the unsorted top as a type, by W9a. For full application
congruence, endpoint generation gives function and argument judgments.
Their endpoint inductions retain the same function product and argument
domain. Apply them; raw argument equality and substitution give equality
of the new and old codomain instances, permitting conversion to \(U\).
Both directions use the same steps with endpoints interchanged.
This covers the full system's simultaneous function/argument replacement,
not only the restricted one-child congruence rules.
For reverse lambda transport, analogously carry the second endpoint's
generated body judgment and body-type sort from the new domain back to
the old domain before applying the reverse body equality implication.
Then reconstruct the first lambda and restore the given type. Compatible
raw beta equality of each mapped domain/body or function/argument pair
gives the asserted raw lambda, product, or application equality in
these congruence cases.

For the full beta equality generator, its source body/argument premises
are mapped by the joint induction. The same type-correctness/product
construction as in abstraction gives a CC-typed root redex. Substitution
types its contractum at the same codomain instance. Raw equality is the
mapped root beta step. For endpoint transport from any additionally
given root type \(U\), use CC subject reduction. In the reverse
direction, uniqueness between the contractum's principal instance and
its given type \(U\) gives beta equality; convert the already typed
root to \(U\). Here \(U\) is sorted: if it were \(\Box\), its
equality to the explicitly sorted principal result type would contradict
W9a's impossibility of a typable subject converting to that atom.
Thus this case does not assume unsafe general beta expansion.

Finally consider the original typing conversion. Its subject induction
gives \(TM:TP\). CC type correctness sorts \(TP\), with the same
literal-star exception handled by \(\ast:\Box\); no source type
maps to \(\Box\). The equality endpoint induction applied to this
sort judgment gives \(TQ\) at the same sort, while its raw equality
gives \(TP=_\beta TQ\). CC conversion therefore licenses the desired
\(TM:TQ\). This closes the mutual induction, including the otherwise
missing target-type sort premise. No normalization entered transport.

\paragraph{W12d. The actual normalization transfer and its precise scope.}
Each full source-typed object and each valid source logical context has
a CC-typed named image by W12c. Every raw compatible source beta
contraction preserves scope and maps to one target contraction by W12b,
at root, in application children, or in either annotation/body field of
either binder. Hence an infinite source path would map to an infinite
CC path. W9 excludes that path. The full concrete/binding-semantics
1988 calculus, with corrected indexed relocation, is therefore strongly
normalizing for full annotated beta, in every valid open environment.
The restricted logical-reduction generators of p.~104, if considered
with their reflexivity and relocation corrected, map to nonempty finite
beta paths as well: root/context cases map to one step and its explicit
transitivity rule concatenates paths. They are not equated to arbitrary
full beta reduction or to the literal erroneous restricted rule system.

This is a proved normalization-preserving forward bridge, which suffices
for the legacy normalization property. It is not advertised as a
bijection of raw syntax, an inverse to the extra target \(\Box\),
or an equivalence of restricted and full source equality judgments.
Nor does it use the paper's stripped program \(v(M)\) to infer
termination of erased annotations. The original p.~110 Corollary~1
proves SN for a closed stripped image; p.~105 refers full constructed-term
normalization to [15], and p.~112 uses it for a normal-type presentation.
W9 and this bridge supply the separate annotated argument editorially.
Primitive recursion, constants-as-extra-rewrite-generators, the optional
universe hierarchy, and Type:Type in section~8 are not in the base claim.
No eta generator occurs in the read base conversion rules. The immutable
1988 source remains unchanged; adoption of its explicit corrective
presentation and routing of a P7 reader result are lead-owned work.
